\chapter[Resultados Preliminares]{Resultados Preliminares}
\label{cap:resultados}

Este capítulo relata o que foi produzido até o fim do TCC1, organizado pelas
atividades da DSR (Tabela~\ref{tab:dsr}).

\section{Identificação do problema}

O diagnóstico da plataforma Gov Hub Cidades mostrou que:

\begin{itemize}
	\item o repositório não continha nenhum componente de ML: não havia
	código de treino, registro de experimentos, datasets versionados nem
	monitoramento de modelos;
	\item a identificação de obras em risco dependia de uma regra de limiares
	fixos que nunca havia sido confrontada com as paralisações observadas;
	\item os dados necessários para um modelo preditivo já existiam nas
	camadas prata, com histórico mensal de dezembro de 2019 a julho de 2026.
\end{itemize}

\section{Definição dos objetivos da solução}

Os requisitos R1--R8 (Tabela~\ref{tab:requisitos}) foram definidos a partir
da literatura e das regras de engenharia da plataforma. A decisão de estender
a plataforma existente, em vez de construir uma plataforma de ML separada, foi
registrada no ADR 0001 do repositório do projeto, com status ``proposto'', em
27 de outubro de 2026.

\pendente{Atualizar o status do ADR 0001 (proposto/aceito) e citar o ADR
0002, sobre o MLflow, quando for escrito.}

\section{Projeto e desenvolvimento}

Até esta etapa foram produzidos:

\begin{enumerate}
	\item a estrutura de pastas do ML no repositório (\texttt{ml/} e
	\texttt{ml/prevencao\_paralisacao/}), com a documentação do ciclo de vida;
	\item a especificação do problema (fase 1 do CRISP-DM), com pergunta,
	unidade de análise, escopo, referência de comparação e decisões em aberto
	(Seção~\ref{sec:caso});
	\item o levantamento inicial dos dados, cujos achados principais (cobertura,
	trocas de fonte, volume de transições) estão na Seção~\ref{sec:caso}.
\end{enumerate}

\section{Análise exploratória dos dados}

A análise exploratória (fase 2 do CRISP-DM) lê as tabelas prata como o
pipeline as monta hoje, sem criar novas tabelas. Ela segue seis perguntas, na
ordem:

\begin{enumerate}
	\item Qual a cobertura e o grão por modalidade e por fonte, e onde há
	lacunas na série mensal?
	\item Quais transições de situação são artefatos da troca de fonte em 2024,
	e como tratá-las?
	\item Como mapear os códigos de situação da planilha mensal e do FGTS para a
	situação canônica?
	\item Qual a taxa de paralisação por modalidade, UF, agente financeiro e
	construtora?
	\item Qual o tempo até a paralisação e até a retomada (análise de
	sobrevivência)?
	\item Quais sinais antecedem a paralisação: ritmo da execução física,
	diferença físico-financeira, tempo sem liberação de recursos, tempo sem
	vistoria?
\end{enumerate}

\pendente{Inserir os resultados da EDA à medida que os notebooks forem
concluídos. Sugestão de figuras: (a) número de obras por mês e por fonte, com
as trocas de fonte marcadas; (b) taxa mensal de transição
\texttt{em\_obras}$\rightarrow$\texttt{paralisada} por fonte; (c) curva de
Kaplan--Meier do tempo até a paralisação por modalidade; (d) distribuição da
diferença físico-financeira em obras que paralisaram e que não paralisaram.}

\pendente{Calcular o desempenho da regra de gargalo vigente
(\textit{recall@k} e \textit{precision@k}) no período SFTP. Este é o primeiro
número que o modelo terá de superar e já pode ser obtido sem nenhum modelo.}

\section{Ameaças à validade}

\begin{itemize}
	\item \textbf{Validade de construto}: o rótulo ``paralisada'' depende da
	fonte, e as trocas de fonte produzem transições que podem não corresponder a
	paralisações reais. A mitigação é restringir o rótulo ao período SFTP e
	tratar a troca de fonte na preparação dos dados;
	\item \textbf{Validade interna}: features calculadas com informação futura
	inflariam artificialmente o desempenho. A mitigação é o teste automatizado
	de vazamento (AV4) e a divisão temporal;
	\item \textbf{Validade externa}: a arquitetura é demonstrada em uma única
	plataforma e em um único caso. A generalização para outros órgãos do Gov Hub
	será discutida a partir do requisito R8 (adição de um segundo modelo sem
	alterar o primeiro), mas não será demonstrada empiricamente;
	\item \textbf{Viés do pesquisador}: os autores participam do desenvolvimento da
	plataforma, o que é inerente à DSR. A mitigação é registrar as decisões em
	ADRs e basear a avaliação em critérios externos (P1--P9 e o \textit{ML Test
	Score}).
\end{itemize}
